\documentclass[11pt]{article}
\usepackage[T1]{fontenc}
\usepackage{lmodern}
\usepackage{amsmath,amssymb,amsthm}
\usepackage[margin=1in]{geometry}
\usepackage{microtype}
\usepackage[hidelinks]{hyperref}

\newtheorem{theorem}{Theorem}
\newtheorem{lemma}[theorem]{Lemma}
\newtheorem{remark}[theorem]{Remark}
\newcommand{\F}{\mathbb F}
\newcommand{\PG}{\operatorname{PG}}
\newcommand{\GL}{\operatorname{GL}}

\title{An Elementary Difference-Spectrum Certificate for\\
Projective Equivalence without CCZ Equivalence over $\F_{11}$}
\author{Xiaojian Zeng\\
PBC School of Finance, Tsinghua University, 100190\\
\texttt{xiaojianzeng@foxmail.com}}
\date{}

\begin{document}
\maketitle

\begin{abstract}
For a function $h\colon\F_{11}\to\F_{11}$, let $\Gamma_h$ be its affine
graph and let $S_h$ be the corresponding $12$-point projective set obtained
by adjoining $(0:1:0)$.  We give a direct exact certificate that
$f(x)=x^4$ and $g(x)=x^8$ are projectively equivalent but not
CCZ-equivalent.  Swapping the second and third projective coordinates maps
$S_f$ onto $S_g$, including the zero and infinite points.  On the affine
side, the complete ordered-difference multisets have histograms
$\{0:50,1:50,3:20\}$ and $\{0:30,1:70,2:20\}$, respectively, so no
invertible affine map can carry one graph to the other.  A general
classification theorem of Dempwolff already implies the same CCZ
inequivalence for these two power functions; the purpose here is the short,
self-contained difference-spectrum proof and its reproducible finite
certificate, not a priority claim.
\end{abstract}

\section{Context and statement}

All arithmetic below is in the prime field $\F_{11}$.  For
$h\colon\F_{11}\to\F_{11}$, set
\[
  \Gamma_h=\{(x,h(x)):x\in\F_{11}\}\subset\F_{11}^2
\]
and
\[
  S_h=\{(x:h(x):1):x\in\F_{11}\}\cup\{(0:1:0)\}
       \subset\PG(2,11).
\]
Huczynska, Klawuhn and Paterson study the projective comparison at
issue here, prove in their Example~2.11 that $x^4$ and $x^8$ are
projectively equivalent, and observe that their different degrees make
them EA-inequivalent~\cite{HuczynskaKlawuhnPaterson2026}.  They define CCZ
equivalence by affine equivalence of graphs, following the equivalence
introduced by Carlet, Charpin and Zinoviev~\cite{CarletCharpinZinoviev1998},
and formulate the existence of a projectively equivalent but
CCZ-inequivalent pair as their Open Problem~4.

There is an important prior-literature qualification.  Dempwolff's general
classification says that the power functions $x^k$ and $x^\ell$ on
$\F_{p^n}$ are CCZ-equivalent precisely when, for some $0\leq a<n$,
\[
 \ell\equiv p^a k\pmod{p^n-1}
 \quad\hbox{or}\quad
 k\ell\equiv p^a\pmod{p^n-1}.
\]
For $p=11,n=1,k=4,\ell=8$, neither condition holds: $8\not\equiv4$
and $4\cdot8\equiv2\not\equiv1\pmod {10}$.  Thus that classification
already gives the non-CCZ conclusion~\cite{Dempwolff2018}.  A later
correction to an intermediate lemma explicitly retains the main
classification theorem~\cite{Dempwolff2022Correction}.  The result below
records instead a direct certificate requiring no classification machinery.

\begin{theorem}\label{thm:main}
Let $f(x)=x^4$ and $g(x)=x^8$ on $\F_{11}$.  The projectivity induced by
\[
 M=\begin{pmatrix}
       1&0&0\\
       0&0&1\\
       0&1&0
    \end{pmatrix}
\]
maps $S_f$ onto $S_g$.  No invertible affine transformation of
$\F_{11}^2$ maps $\Gamma_f$ onto $\Gamma_g$.  Consequently $f$ and $g$
are projectively equivalent but not CCZ-equivalent.
\end{theorem}

\section{The projective certificate}

The determinant of $M$ is $-1=10$ in $\F_{11}$, so $M$ induces a
projectivity.  It exchanges
\[
 M(0:0:1)=(0:1:0),\qquad M(0:1:0)=(0:0:1),
\]
thereby accounting for both exceptional points.  If $x\ne0$, then
\[
 M(x:x^4:1)=(x:1:x^4)=(x^{-3}:x^{-4}:1).
\]
Put $u=x^{-3}$.  Since $\gcd(3,10)=1$, the map $x\mapsto u$ permutes
$\F_{11}^{*}$; moreover
\[
 u^8=x^{-24}=x^{-4},
\]
because exponents of nonzero elements may be reduced modulo $10$.
It follows that the ten remaining images are exactly
$(u:u^8:1)$ for $u\in\F_{11}^{*}$.  Hence $M(S_f)=S_g$.

\section{The affine-difference obstruction}

For a subset $\Gamma\subset\F_{11}^2$ and a nonzero vector $w$, define
the ordered-difference count
\[
 \Delta_\Gamma(w)
   =\#\{(P,Q)\in\Gamma^2:P-Q=w\}.
\]
The differential spectrum is standard machinery in the study of CCZ
equivalence~\cite{Kaleyski2021}.  We include the short invariance proof,
which applies directly in odd characteristic and also incorporates vertical
differences.

\begin{lemma}\label{lem:invariant}
The multiset
\[
 \mathcal D(\Gamma)=
 \bigl\{\!\bigl\{\Delta_\Gamma(w):
       w\in\F_{11}^2\setminus\{0\}\bigr\}\!\bigr\}
\]
is invariant under invertible affine transformations of $\F_{11}^2$.
\end{lemma}

\begin{proof}
Write an affine permutation as $T(z)=Az+t$, where
$A\in\GL(2,11)$.  For every ordered pair,
\[
 T(P)-T(Q)=A(P-Q).
\]
The map $(P,Q)\mapsto(T(P),T(Q))$ is a bijection, and $A$ permutes the
$120$ nonzero vectors.  Therefore
$\Delta_{T(\Gamma)}(Aw)=\Delta_\Gamma(w)$ for every $w\ne0$, proving the
claim.
\end{proof}

For $d\in\{4,8\}$ abbreviate
$\Delta_d=\Delta_{\Gamma_{x^d}}$.  If $a\ne0$, an ordered difference
$(a,b)$ has the form
\[
 (y+a,(y+a)^d)-(y,y^d).
\]
After writing $y=at$, we obtain the reconstruction formula
\begin{equation}\label{eq:reconstruction}
 \Delta_d(a,b)=
 \#\left\{t\in\F_{11}:(t+1)^d-t^d=b/a^d\right\}.
\end{equation}
For $a=0$ and $b\ne0$, the count is zero.  Thus the two rows below,
computed directly modulo $11$, reconstruct every one of the $120$ entries.

\begin{center}
{\small
\begin{tabular}{c|rrrrrrrrrrr}
$d\backslash t$&0&1&2&3&4&5&6&7&8&9&10\\ \hline
$4$&1&4&10&10&6&0&5&1&1&7&10\\
$8$&1&2&2&4&6&0&5&7&9&9&10
\end{tabular}
\par\medskip
\begin{tabular}{c|rrrrrrrrrrr}
$d\backslash c$&0&1&2&3&4&5&6&7&8&9&10\\ \hline
$4$&1&3&0&0&1&1&1&1&0&0&3\\
$8$&1&1&2&0&1&1&1&1&0&2&1
\end{tabular}
}
\end{center}

The first table lists $(t+1)^d-t^d$ for $t=0,\ldots,10$; the second
lists the multiplicity of each value $c=0,\ldots,10$.  For each of the ten
choices $a\ne0$, multiplication by $a^{-d}$ merely permutes the eleven
values of $b$ in~\eqref{eq:reconstruction}.  Repeating the second row ten
times and adjoining the ten forced vertical zero entries gives
\[
 \mathcal D(\Gamma_f)=\{0:50,\ 1:50,\ 3:20\},
 \qquad
 \mathcal D(\Gamma_g)=\{0:30,\ 1:70,\ 2:20\}.
\]
Here, for example, $\{0:50,1:50,3:20\}$ means that the values $0,1,3$
occur with multiplicities $50,50,20$.  Each multiset has $120$ entries,
and the weighted sum of its entries is $110$, the number of ordered pairs
of distinct graph points.  Since the two multisets differ, Lemma
\ref{lem:invariant} rules out every affine graph equivalence.  Together
with the projective certificate, this proves Theorem~\ref{thm:main}.

\section{Reproducibility and assistance disclosure}

The accompanying file \path{research_exact_certificate.json} contains
all twelve projective images and both complete $120$-entry difference
tables.  The standard-library Python checker
\path{research_independent_check.py} reconstructs the tables from
\eqref{eq:reconstruction}, recomputes them independently by looping over
ordered graph pairs, and requires entrywise agreement.  All operations are
integer operations reduced modulo $11$; there is no floating-point or
randomized computation.  The displayed derivation proves the theorem, while
the files provide reproducible corroboration.  This distinction between an
algebraic proof and an exact-arithmetic companion verifier is also made in
recent prime-field work of Zeng, Liu and Ratnavelu
\cite{ZengLiuRatnavelu2026}; that preprint is cited here as the most closely
related verification practice in the author's supplied bibliography, not
as a source for the present CCZ argument.

OpenAI Codex assisted with exact finite-field computation, proof
organization, literature discovery and LaTeX preparation.  A separate
automated audit recomputed the finite certificate and checked the displayed
argument; such an audit is not independent peer review.  All theorem-level
claims used in the paper are stated and proved above, or are attributed to
the cited primary literature.

\bibliographystyle{abbrv}
\bibliography{references}

\end{document}
